\documentclass[11pt,a4paper,oneside]{report}             % Single-side
%\documentclass[11pt,a4paper,twoside,openright]{report}  % Duplex

%\PassOptionsToPackage{chapternumber=Huordinal}{magyar.ldf}
\usepackage{t1enc}
\usepackage[latin2]{inputenc}
\usepackage[magyar, german, english]{babel}
\usepackage{mathtools}
\usepackage{amssymb}
\usepackage{siunitx}
\usepackage{rotating}
\usepackage{enumerate}
\usepackage{pgfplots, pgfplotstable}
\usepackage[thmmarks]{ntheorem}
\usepackage{graphics}
\usepackage[refpage]{nomencl}
\usepackage{glossaries}
\usepackage{multirow}
\usepackage{graphicx,floatpag,fancyhdr}
\usepackage{epsfig}
\usepackage{listings}
\usepackage{makecell}
\usepackage{booktabs}
\usepackage{color}
\usepackage{lastpage}
\usepackage{anysize}
\usepackage{sectsty}
\usepackage{array}
\usepackage{setspace}  % Ettol a tablazatok, abrak, labjegyzetek maradnak 1-es sorkozzel!
\usepackage[hang]{caption} 
\usepackage{subcaption}
\usepackage{hyperref}
\usepackage{afterpage}
\usepackage{booktabs}
% Own inclusions
\usepackage{titlesec}
\usepackage{relsize}
\usepackage{anyfontsize}
\usepackage{pdfpages}
\usepackage{placeins}
\usepackage{makeidx}
%\usepackage{adjustbox}
\usepackage[export]{adjustbox}
\usepackage{enumitem}
\usepackage[inkscapeopt= -D -z]{svg}

%--------------------------------------------------------------------------------------
%	SVG package setup for an easier life ;)
%--------------------------------------------------------------------------------------
\setsvg{inkscape={"C:/Program Files/Inkscape/inkscape" -D -z}}

%--------------------------------------------------------------------------------------
%	Path setups: figures, svgs, tables, charts
%--------------------------------------------------------------------------------------
\svgpath{{"C:/Users/{"Barancsuk_Lilla"}/Documents/{"current_semester"}/Dipterv_repo/{"Latex_document"}/figures/"}}
\graphicspath{{figures/}{figures/for_timelines/}{figures/for_frames/}{figures/GT/}{figures/calcorr/fovam/}{figures/calcorr/ibike/}{figures/calcorr/}{figures/calcorr/uncalib/}{figures/errors/}{figures/appendix/advert_app/}{figures/appendix/timelines/}{figures/appendix/video_mixer/}{figures/error_timelines/}}

\newcommand{\tablepath}{tables/}

\makeatletter
\newcommand{\includetable}[1]{%
	\@ifundefined{tablepath}{%
		\InputIfFileExists{#1}{}{}%
	}{%
		\InputIfFileExists{\tablepath/#1}{}{\InputIfFileExists{#1}{}{}}%
	}
}
\makeatother 

\newcommand{\chartpath}{charts/}

\makeatletter
\newcommand{\includechart}[1]{%
	\@ifundefined{chartpath}{%
		\InputIfFileExists{#1}{}{}%
	}{%
		\InputIfFileExists{\chartpath/#1}{}{\InputIfFileExists{#1}{}{}}%
	}
}
\makeatother 

%--------------------------------------------------------------------------------------
%	Main variables
%--------------------------------------------------------------------------------------
\newcommand{\vikszerzo}{Lilla \textsc{Barancsuk}}
\newcommand{\vikHunszerzo}{Barancsuk Lilla}
\newcommand{\vikkonzulens}{Krist�f \textsc{Csorba}, PhD}
\newcommand{\vikcim}{Development and integration of camera based traffic monitoring system}
\newcommand{\viktanszek}{Department of Automation and Applied Informatics}
\newcommand{\vikdoktipus}{Masters thesis}
\newcommand{\vikdepartmentr}{Barancsuk Lilla}

%--------------------------------------------------------------------------------------
%	Page layout setup
%--------------------------------------------------------------------------------------
% we need to redefine the pagestyle plain
% another possibility is to use the body of this command without \fancypagestyle
% and use \pagestyle{fancy} but in that case the special pages
% (like the ToC, the References, and the Chapter pages)remain in plane style

\pagestyle{plain}
%\setlength{\parindent}{0pt} % �ttekinthet�bb, angol nyelv� dokumentumokban jellemz�
%\setlength{\parskip}{8pt plus 3pt minus 3pt} % �ttekinthet�bb, angol nyelv� dokumentumokban jellemz�
\setlength{\parindent}{12pt} % magyar nyelv� dokumentumokban jellemz�
\setlength{\parskip}{0pt}    % magyar nyelv� dokumentumokban jellemz�

\marginsize{35mm}{25mm}{15mm}{15mm} % anysize package
\setcounter{secnumdepth}{0}
\sectionfont{\large\upshape\bfseries}
\setcounter{secnumdepth}{2}
\singlespacing
\frenchspacing

% This defines the fancy page style to be similar to plain
\pagestyle{fancy}
\fancyhf{}% Clear header/footer
\fancyfoot[C]{\thepage}
\renewcommand{\headrulewidth}{0pt}% Remove header rule

% Page style plainlower is similar to plain, but lowers page number by 6 lines of text
\fancypagestyle{plainlower}{
	\fancyhf{}% Clear header/footer
	\fancyfoot[C]{\raisebox{-6\baselineskip}{\thepage}}
	\renewcommand{\headrulewidth}{0pt}% Remove header rule
}

%--------------------------------------------------------------------------------------
%	Table settings
%--------------------------------------------------------------------------------------
\renewcommand\theadalign{cb}
\renewcommand\theadfont{\bfseries}
\renewcommand\theadgape{\Gape[4pt]}
\renewcommand\cellgape{\Gape[4pt]}
\renewcommand{\cellalign}{cc}
\newcolumntype{?}{!{\vrule width 1pt}}

%--------------------------------------------------------------------------------------
%	hmmm???
%--------------------------------------------------------------------------------------
%\parfillskip 0pt plus 0.75\textwidth
%\widowpenalty=10000
%\clubpenalty=10000

%--------------------------------------------------------------------------------------
%	Setup hyperref package
%--------------------------------------------------------------------------------------
\hypersetup{
    %bookmarks=true,            % show bookmarks bar?
    unicode=false,             % non-Latin characters in Acrobat�s bookmarks
    pdftitle={\vikcim},        % title
    pdfauthor={\vikszerzo},    % author
    pdfsubject={\vikdoktipus}, % subject of the document
    pdfcreator={\vikszerzo},   % creator of the document
    pdfproducer={Producer},    % producer of the document
    pdfkeywords={keywords},    % list of keywords
    pdfnewwindow=true,         % links in new window
    colorlinks=true,           % false: boxed links; true: colored links
    linkcolor=black,           % color of internal links
    citecolor=black,           % color of links to bibliography
    filecolor=black,           % color of file links
    urlcolor=black             % color of external links
}

%--------------------------------------------------------------------------------------
%	Set up listings
%--------------------------------------------------------------------------------------
\lstset{
	basicstyle=\scriptsize\ttfamily, % print whole listing small
	keywordstyle=\color{black}\bfseries\underbar, % underlined bold black keywords
	identifierstyle=, 					% nothing happens
	commentstyle=\color{white}, % white comments
	stringstyle=\scriptsize\sffamily, 			% typewriter type for strings
	showstringspaces=false,     % no special string spaces
	aboveskip=3pt,
	belowskip=3pt,
	columns=fixed,
	backgroundcolor=\color{lightgray},
} 		
\def\lstlistingname{code-listing}	

%--------------------------------------------------------------------------------------
%	Some new commands and declarations
%--------------------------------------------------------------------------------------
\newcommand{\code}[1]{{\upshape\ttfamily\scriptsize\indent #1}}
\newcommand{\reg}{\textsuperscript{\textcircled{\textsc r}}}
% define references
\newcommand{\figref}[1]{\ref{fig:#1}.}
\renewcommand{\eqref}[1]{(\ref{eq:#1})}
\newcommand{\listref}[1]{\ref{listing:#1}.}
\newcommand{\sectref}[1]{\ref{sect:#1}}
\newcommand{\tabref}[1]{\ref{tab:#1}.}
\newcommand{\hsp}{\hspace{0pt}}

%--------------------------------------------------------------------------------------
%	Math operators
%--------------------------------------------------------------------------------------
\DeclareMathOperator*{\argmax}{arg\,max}
%\DeclareMathOperator*[1]{\floor}{arg\,max}
\DeclareMathOperator{\sign}{sgn}
\DeclareMathOperator{\rot}{rot}
\DeclareMathOperator{\avg}{avg}
\DeclareMathOperator{\avgt}{\underset{\textit{t}}{avg}}
%Math mode set\part{title}
%\everymath{\displaystyle}

%--------------------------------------------------------------------------------------
%	Color declarations
%--------------------------------------------------------------------------------------
\definecolor{lightgray}{rgb}{0.95,0.95,0.95}
\definecolor{indigo}{RGB}{47, 27, 60}
\definecolor{darkviolet}{RGB}{62, 44, 102}
\definecolor{lightviolet}{RGB}{96, 92, 153}
\definecolor{taupe}{RGB}{176, 158, 141}
\definecolor{lighttaupe}{RGB}{216, 208, 197}
\definecolor{camel}{RGB}{232, 208, 174}

%--------------------------------------------------------------------------------------
%	Chart color cycle
%--------------------------------------------------------------------------------------
\pgfplotscreateplotcyclelist{violets}{
	{fill=indigo!60!white, draw=indigo},
	{fill=taupe!60!white, draw=taupe},
	{fill=darkviolet!60!white, draw=darkviolet},
	{fill=camel!60!white, draw=camel},
	{fill=lightviolet!60!white, draw=lightviolet},
}

%--------------------------------------------------------------------------------------
%	Chart setup
%--------------------------------------------------------------------------------------
\newcommand{\placelegendWithRequirement}{
	legend style={
		at={(1.28,1)}, 
		anchor=north, 
		legend columns=1}, 
	reverse legend,
	legend cell align=left,}	

\newcommand{\placelegend}{
	legend style={
		at={(1.2,1)}, 
		anchor=north, 
		legend columns=1}, 
	reverse legend,
	legend cell align=left,}	
\newcommand{\custombarwidth}{0.13cm}

%--------------------------------------------------------------------------------------
%	Includes
%--------------------------------------------------------------------------------------
\author{\vikszerzo}
\title{\viktitle}
\includeonly{
	guideline,%
	project,%
	titlepage,%
	declaration,%
	abstract,%
	introduction,%
	chapter1,%
	%chapter2,%
	chapter3,%
	chapter4,%
	chapter5,%
	conclusions,%
	acknowledgement,%
	appendices,%
}

%--------------------------------------------------------------------------------------
%	Setup captions
%--------------------------------------------------------------------------------------
\captionsetup[figure]{
%labelsep=none,
%font={footnotesize,it},
%justification=justified,
width=.75\textwidth,
aboveskip=10pt}

\renewcommand{\captionlabelfont}{\small\bf}
\renewcommand{\captionfont}{\footnotesize\it}

%---------------
%  makeindex
%---------------
\makeindex

%-------------------
%  makenomenclature
%-------------------
\makenomenclature
\makeglossaries
%--------------------------------------------------------------------------------------
%	Table of contents and the main text
%--------------------------------------------------------------------------------------
\begin{document}

%-------------------------------------------------------------------
%	Chapter header format with HUGE numbering <3 <3 <3 (Because it's pretty)
%-------------------------------------------------------------------
\titleformat{\chapter}[display]{\flushright\fontseries{b}\fontsize{120}{120}\selectfont}{\fontseries{b}\fontsize{180}{4}\selectfont\thechapter\hsp}{0pt}{\Huge\bfseries}[]

\singlespacing
\pagenumbering{arabic}
\onehalfspacing
\includepdf[pages={1}]{../Kameras-forgalomszamlalo-rendszer-fejlesztese-Feladatkiiras-2.pdf}
%--------------------------------------------------------------------------------------
%	The title page
%--------------------------------------------------------------------------------------
\begin{titlepage}
\begin{center}
\includegraphics[width=60mm,keepaspectratio]{figures/BMElogo.png}\\
\vspace{0.3cm}
\textbf{Budapest University of Technology and Economics}\\
\textmd{Faculty of Electrical Engineering and Informatics}\\
\textmd{\viktanszek}\\[5cm]

\vspace{0.4cm}
{\huge \bfseries \vikcim}\\[0.8cm]
\vspace{0.5cm}
\textsc{\Large \vikdoktipus}\\[4cm]

\begin{tabular}{cc}
 \makebox[7cm]{\emph{Author}} & \makebox[7cm]{\emph{Advisor}} \\
 \makebox[7cm]{\vikszerzo} & \makebox[7cm]{\vikkonzulens}
\end{tabular}

\vfill
{\large \today}
\end{center}
\end{titlepage}



\tableofcontents\vfill
\listoffigures\addcontentsline{toc}{chapter}{List of figures}
\printglossaries\addcontentsline{toc}{chapter}{List of abbreviations}
%--------------------------------------------------------------------------------------
% Nyilatkozat
%--------------------------------------------------------------------------------------
\begin{center}
\large
\textbf{HALLGAT�I NYILATKOZAT}\\
\end{center}

Alul�rott \emph{\vikHunszerzo}, szigorl� hallgat� kijelentem, hogy ezt a diplomatervet meg nem engedett seg�ts�g n�lk�l, saj�t magam k�sz�tettem, csak a megadott forr�sokat (szakirodalom, eszk�z�k stb.) haszn�ltam fel. Minden olyan r�szt, melyet sz� szerint, vagy azonos �rtelemben, de �tfogalmazva m�s forr�sb�l �tvettem, egy�rtelm�en, a forr�s megad�s�val megjel�ltem.

Hozz�j�rulok, hogy a jelen munk�m alapadatait (szerz�(k), c�m, angol �s magyar nyelv� tartalmi kivonat, k�sz�t�s �ve, konzulens neve) a BME VIK nyilv�nosan hozz�f�rhet� elektronikus form�ban, a munka teljes sz�veg�t pedig az egyetem bels� h�l�zat�n kereszt�l (vagy autentik�lt felhaszn�l�k sz�m�ra) k�zz�tegye. Kijelentem, hogy a beny�jtott munka �s annak elektronikus verzi�ja megegyezik. D�k�ni enged�llyel titkos�tott diplomatervek eset�n a dolgozat sz�vege csak 3 �v eltelte ut�n v�lik hozz�f�rhet�v�.

\begin{flushleft}
\vspace*{1cm}
Budapest, \today
\end{flushleft}

\begin{flushright}
 \vspace*{1cm}
 \makebox[7cm]{\rule{6cm}{.4pt}}\\
 \makebox[7cm]{\emph{\vikHunszerzo}}\\
 \makebox[7cm]{hallgat�}
\end{flushright}
\thispagestyle{empty}

\vfill
\clearpage
\thispagestyle{empty} % an empty page


%----------------------------------------------------------------------------
% Abstract in hungarian
%----------------------------------------------------------------------------
\chapter*{Kivonat}\addcontentsline{toc}{chapter}{Kivonat}

Amellett, hogy a k�zutak forgalmi adatainak nyilv�ntart�sa az �nkorm�nyzatok k�telez? feladata, a forgalomsz�ml�l�s sokf�le �rt�kes inform�ci�t szolg�ltat. Seg�ts�g�vel megbecs�lhet? az adott �tszakasz k�rosanyag-kibocs�t�sa �s terhelts�ge, amivel hozz�j�rul a megl�v? �th�l�zat �zemeltet�s�hez, beruh�z�sok tervez�s�hez. A real-time forgalomsz�ml�l� rendszerek pedig alkalmasak forgalmi dug�k detekt�l�s�ra, megel?z�s�re is.

Napjainkban a k�zi forgalomsz�ml�l�st egyre ink�bb felv�ltj�k az automatiz�lt m�dszerek. 
A radaros �s az �tba �p�thet? (indukt�v hurkos) rendszerek mellett a videofelv�tel alapj�n t�rt�n? sz�ml�l�s �rvend a legnagyobb n�pszer?s�gnek. 
A korl�tozott er?forr�s�, be�gyazott k�rnyezetben azonban tov�bbra is kih�v�st jelent a k�pfeldolgoz�son alapul� sz�ml�l�s val�sidej? megval�s�t�sa. 
Az elterjedt, nagy sz�m�t�sig�ny? h�tt�r-elt�vol�t�st �s feature-pont k�vet�st haszn�l� megold�sok sokszor nem �rik el a k�v�nt pontoss�got vagy sebess�get.

Dolgozatomban egy, a forgalomt�meg becsl�s�re alkalmas rendszert mutatok be, amely az utat fel�lr?l r�gz�t? kamera k�p�t dolgozza fel. 
C�lom egy olyan l�mpatestbe �p�thet? intelligens szenzor megalkot�sa volt, amely adott korl�tos er?forr�s� hardveren is k�pes val�s id?ben m?k�dni, az elhalad� j�rm?vek sz�m�t, t�pus�t adott pontoss�ggal meg�llap�tani, valamint megbecs�li sebess�g�ket �s k�vet�si t�vols�gukat.

Dolgozatomban r�szletesen bemutatom a rendszer magj�t alkot� szoftver m?k�d�s�t �s fel�p�t�s�t, valamint a vide�k�pek feldolgoz�s�nak folyamat�t, melynek alapja t�bbf�le, hat�kony vide�feldolgoz�si m�dszer �sszekapcsol�sa: h�tt�r-elt�vol�t�s, tripwire-alap� sz�ml�l�s, ill. kis sz�m�t�sig�ny? adatgy?jt�s �s oszt�lyoz�s kombin�l�sa, ezek el?nyeinek egyes�t�se.
Kit�rek a rendszert kieg�sz�t?, a val�s forgalomban val� m?k�d�st lehet?v� tev? szoftver- �s hardverk�rnyezetre, a t�volr�l t�rt�n? kalibr�ci�t �s tesztel�st megval�s�t� alkalmaz�sra �s a rendszer konfigur�ci�j�nak menet�re.
R�szletezem az alkalmaz�s pontoss�g- �s se\-bes\-s�g\-tesz\-te\-l�\-s�\-nek mik�ntj�t, a tesztek eredm�ny�t, valamint a tov�bbfejl?d�s lehets�ges �tjait is.
\vfill

%----------------------------------------------------------------------------
% Abstract in english
%----------------------------------------------------------------------------
\chapter*{Abstract}\addcontentsline{toc}{chapter}{Abstract}

The registration of traffic data of public roads is not only a compulsory task of local governments, but the traffic count provides valuable information of various kinds. It can be used to estimate the pollution emission and load of a road, which contributes to the operation of the current facilities and to the appropriate planning of new investments. Moreover, real-time traffic monitoring systems are able to detect and even to avoid traffic jams, too.

Nowadays, manual surveillance are more and more being replaced by automatized methods. Besides the radar-based systems and inductive loop detectors installed in the road, the video-based traffic detection is among the most popular methods. However, its real-time implementation, by means of image processing, is still challenging in resource-limited embedded systems. The popular algorithms of background subtraction and feature point tracking are usually associated with a computational cost being too high to meet the accuracy and speed requirements.

In this work, I introduce a system for traffic flow estimation, based on the processing of a video stream of a camera installed above the road. My goal was to device an intelligent sensor that can be installed within a street-lighting lamp, which operates in real-time on the given resource-limited hardware. It is able to determine the number and type of the vehicles passing by with an acceptable accuracy, and also to estimate their speed and following distance.

The operation and architecture of the core software and the video processing task is presented in detail in this work. The main idea is the combination of various efficient video processing methods: background subtraction, tripwire-based counting, data mining and classification algorithms of low computational cost, respectively. The benefits of each algorithm is exploited via their appropriately combined use.
The supplementary software and hardware elements, which facilitate the remote calibration, testing, and the configuration process are discussed.
Moreover the accuracy and speed testing process and the results are presented along with possible ways of further improvement.
\vfill


%----------------------------------------------------------------------------
\chapter*{Introduction}\addcontentsline{toc}{chapter}{Introduction}
%----------------------------------------------------------------------------
\section{Traffic load estimation}
%----------------------------------------------------------------------------
In recent years, there has been an increased scope for the automatic analysis of urban traffic activity\ref{Buch2011}.
This case is due to the increasing number of vehicles, the need for better understanding traffic dynamics for investment purposes, and also the accessibility of various sensors.
Formerly traffic estimation was performed by operators, by the means of on-the-spot observation or video analysis.
By the automation of these processes the main concept of traffic analysis is to aid or fully avoid these human operators.

The information obtained by traffic monitoring systems are widely used to regulate traffic flow and contribute to administrative decisions about transportation infrastructure development, maintenance and investments \ref{MagyarKozut}.
In addition to the detection and classification of road users, which is an elementary task, several other monitoring objectives can be supported. 
Typical traffic data such as flow, density, and average traffic speed can be extracted.
Some systems also recognize various situations including traffic violations or accidents (e.g., cars, motorbikes, and pedestrians).
Other real-time applications notify users about congestions, roadworks and accidents, or control traffic lights in order to influence traffic flow direction \ref{AzoSensor, Thiruverahan2015}.

!!!! there is now strong interest in environmental factors such as greenhouse gases, pollutant emissions, and fuel consumption. It is now possible to combine high-resolution real-time traffic data with instantaneous emission models to estimate these environmental measures in real time. \ref{Morris2012a}

In the long term surveillance systems will be integral parts of the intelligent city infrastructure, identifying and influencing these city's traffic dynamics\ref{enLight}.

\subsection{Traffic monitoring techniques}
Nowadays a variety of sensing modalities is available for on-road vehicle detection.

Inductive loop systems installed in the road are considered the most reliable traffic classification and detection method available, larger metal objects can be accurately detected this way, on the other hand the installation and maintenance is time-consuming and difficult\ref{Diamond}.

Other robust detection methods including radar \ref{DeepBlue}, infrared \ref{Swarco} and laser \ref{SICK, Gallego2009} sensors are also popular, although a complex machinery needs to be deployed.

Weight-based classification and speed detection is also possible using a piezoelectric sensor\ref{Te, Rivas2017}.

Az ipar számos megoldást kínál erre a problémára. Elterjedtek például az útba telepíthet˝o,
induktív hurkos rendszerek. Ezek pontosan azonosítanak nagyobb fémtárgyakat, ugyanakkor utólagos
beszerelés esetén az úttest felbontását igényik, így karbantartásuk is nehéz [1]. A radaros [2],
infravörös fénnyel [3] vagy lézerrel [4] m˝uköd˝o szenzorok is népszer˝uek, ehhez azonban bonyolult
berendezéseket kell telepíteni. Esetenként a járm˝uvek súlyát használják detektálásra, például
piezo-kristályokon alapuló tömegméréssel – ehhez azonban szintén az úttest megbontására van
szükség [5].

\subsection{Some results in video-based traffic monitoring}
Video cameras have been deployed for a long time for traffic monitoring purposes, because they provide a rich information source for human understanding\ref{Buch2011}.
In addition to the detection and classification of road users, which is an elementary task, some advanced sensing systems can track, measure the velocity, or understand the complex behaviour of the bypassing vehicles.
Automatic license plate recognition is also a extensively used\ref{Sivaraman2013}. 

 Imaging technology has immensely progressed
in recent years. Cameras are cheaper, smaller, and of higher
quality than ever before. Concurrently, computing power has
dramatically increased. Furthermore, in recent years, we have
seen the emergence of computing platforms geared toward
parallelization, such as multicore processing and graphical processing
units (GPUs). Such hardware advances allow computer
vision approaches for vehicle detection to pursue real-time
implementation.
With advances in camera sensing and computational technologies,
advances in vehicle detection using monocular vision,
stereo vision, and sensor fusion with vision have been
an extremely active research area in the intelligent vehicles
community. On-road vehicle tracking has been also extensively
studied. It is now commonplace for research studies to report
the ability to reliably detect and track on-road vehicles in real
time, over extended periods [3], [5].


In urban environments, several monitoring objectives can
be supported by the application of computer vision and pattern
recognition techniques, including the detection of traffic
violations (e.g., illegal turns and one-way streets) and the
identification of road users (e.g., vehicles, motorbikes, and
pedestrians).

Video analytics may now provide added value to cameras by automatically
extracting relevant information. This way, computer
vision and video analytics become increasingly important for
intelligent transport systems (ITSs).

One level up, detected vehicles
are associated across frames, allowing for vehicle tracking.
%----------------------------------------------------------------------------
\section{The SOLSUN project}
%----------------------------------------------------------------------------
\subsection{Project goals and constraints}
%precision, speed constraints, SW, HW

%----------------------------------------------------------------------------
\section{Thesis goals and organization}
%----------------------------------------------------------------------------

%----------------------------------------------------------------------------
\setcounter{chapter}{0}%
\renewcommand{\thechapter}{\arabic{chapter}}%
\chapter{Basic concepts and terms}\label{chap:Concepts}
%----------------------------------------------------------------------------
In the following chapter certain basic definitions and ideas, that are necessary for understanding the operation of the Traffic Sensor, will be discussed and formalized.

Let one pixel of a red-green-blue (RGB) image be denoted by $\boldsymbol{I}(x,y,t)$, where $x, y$ are the variables of the discrete spatial axes, and $t$ notes the discrete time.
$\boldsymbol{I}(x,y,t)$ is a three-element vector consisting of the intensities of each channel of an RGB image.
%----------------------------------------------------------------------------
\begin{figure}[bh]
	\centering
	\includesvg[width=\textwidth,pretex=\Large]{TI_creation.pdf_tex}
	\caption{The generation procedure of a timeline image. Each column corresponds to a tripwire state. The horizontal axis represents time.\label{fig:TI_creation}}
\end{figure}
%----------------------------------------------------------------------------
\begin{figure}[h!]
	\centering
	\scalebox{.7}{\input{figures/notations.pdf_tex}}
	\caption{Notations used in the thesis, and relationship between the discussed data structures. Each frame is an element of a Frame-strip, whilst the tripwire is a straight line on each frame. The consecutive states of the tripwire are considered a timeline image, where the detection takes place.	\label{fig:notations}}
\end{figure}
%----------------------------------------------------------------------------
\section{Frame}
Let a frame at a certain $t$ time be denoted by
\begin{gather*}
	\boldsymbol{F}_t=\boldsymbol{F}_t(x,y)=\boldsymbol{I}(x,y,t),\\
	x=0,1,2,\dotsc,n_{\text{cols}}-1,\\
	y=0,1,2,\dotsc,n_{\text{rows}}-1,
\end{gather*}
where $n_{\text{cols}}$ stands for the number of columns, and $n_{\text{rows}}$ is the number of rows in a frame.
Each RGB frame is a matrix of pixels composed of three-element vectors.
Grayscale images are built up by one-element pixels.

\section{Frame-strip}
%----------------------------------------------------------------------------
A Frame-strip is a sequence of frames in the order of recording, denoted by:
\begin{gather*}
	\boldsymbol{FS}=\boldsymbol{FS}(t)=\boldsymbol{F}_t(x,y), \\
	t=0,1,2,\dotsc,n_{\text{frames}}-1,
\end{gather*}
where $n_{\text{frames}}$ denotes the frame count in a video sequence.
%--------------------------------------------------------------------------
\section{Tripwire}
%----------------------------------------------------------------------------
The tripwire is a straight line on each frame, where the detection takes place (figure \ref{fig:TI_creation}.).
Let the intensities of the tripwire at a certain $t$ time be denoted by 
\begin{displaymath}
\boldsymbol{T}_t=\boldsymbol{F}_t(x,y)=\boldsymbol{T}_t(k),
\end{displaymath}
where $x$ and $y$ form a 4-connected line on the image pane of the $t$th frame, and the series of the successive pixels of the line are numbered with $k$:
\begin{displaymath}
	k=0,1,\dotsc,l_{\text{tripwire}}-1,
\end{displaymath}
with $l_{\text{tripwire}}$ being the number of pixels of the line, that is the length of the tripwire.

For detection purposes the tripwire must be approximately perpendicular to the traffic direction.
To ensure this angle the tripwire is configured manually in the Traffic Sensor application .
In recent scholarly articles the tripwire is mostly called either virtual detection line (VDL) or check-line.
%----------------------------------------------------------------------------
\section{Timeline image}
%----------------------------------------------------------------------------
A timeline image (TI) is a sequence of the the states of the tripwire on successive frames (figure \ref{fig:notations}.), denoted by
\begin{displaymath}
	\boldsymbol{TI}=\boldsymbol{TI}(t,k),
\end{displaymath}
where the horizontal axis of the TI represents time ($t$), and the vertical axis corresponds to the pixels of the tripwire itself ($k$).
 
In other words, the $i$th column in the timeline image is filled with the tripwire pixels of the $i$th time frame:
\begin{displaymath}
	\boldsymbol{TI}(t=i,k)=\boldsymbol{T}_i,
\end{displaymath}
and the $\text{i}$th row corresponds to the $\text{i}$th tripwire pixel changes throughout the video stream:
\begin{displaymath}
\boldsymbol{TI}(t,k=i)=\boldsymbol{T}(i),
\end{displaymath}

The generation process of a TI is depicted in figure \ref{fig:TI_creation}.
In the literature timeline images are mostly mentioned as stack of lines or time-spatial images.
%----------------------------------------------------------------------------

%----------------------------------------------------------------------------
\chapter{The core software}\label{chap:Core software}
%----------------------------------------------------------------------------
This chapter presents the structure and operation procedure of the core software system of the Traffic Sensor, in particular with regard to the effective image processing algorithms and software development paradigms, that facilitate the real-time operation of the system.

The core software is responsible for the processing of the video stream itself.
Besides that other 
%----------------------------------------------------------------------------
\section{Software architecture}
%----------------------------------------------------------------------------
In this section the architectural parts of the framework, including the data storage technique, the processing method and the parallel operation of the Traffic Sensor core software are discussed.

\subsection{Media storage}
Timeline images can either be extracted from transformed versions of the original frame, or created from the latter calculated data.

\subsubsection{Timeline images}
\begin{enumerate}
	\item \textbf{Original Timeline}
	\item \textbf{MOG2 Timeline} 
	\item \textbf{Morphology Timeline}
	\item \textbf{Motionless Timeline}
	\item \textbf{Size, Speed, Following Distance Timeline: Parameter Timelines}
\end{enumerate}

\subsubsection{Frame-strips}
\begin{enumerate}
	\item \textbf{Raw Frame-strip}
	\item \textbf{Original Frame-strip}
	\item \textbf{MOG2 Frame-strip}
\end{enumerate}

\section{Plug-in architecture, Processors}
Processors

\section{Parallel processing}
%----------------------------------------------------------------------------
%----------------------------------------------------------------------------
\section{Processing steps}
%----------------------------------------------------------------------------
The processing of the video-stream consists of four main stages, as depicted in figure \ref{fig:processing_steps}.

\begin{figure}[bh]
	\centering
	\scalebox{0.7}{	\LARGE\input{figures/full_system_flowchart.pdf_tex}}
	\caption{Steps of the processing of the video stream in the core software of the Traffic Sensor.\label{fig:processing_steps}}
\end{figure}

The first step is preprocessing, that is the remapping of each frame to achieve a standard form for effective processing.
Second, a background-subtraction and its post-corrections are applied to detect the moving vehicles precisely.
The third stage is data extraction, including vehicle detection, classification and parameter calculation.
The final step is evaluation and testing. 
This step is optional, and is only available when the system is tested during development and the operation is not continuous.

In each processing stage a series of different timeline images and frame-strips are used.
%----------------------------------------------------------------------------
\subsection{Preprocessing}
%----------------------------------------------------------------------------
In the preprocessing stage a series of transforms are performed on each Raw Frame as seen on figure \ref{fig:transforms}.
Remapping creates a standard frame-scheme regardless of the video file format, camera type, placement and environment.
This regular form has a pre-defined size, that is small enough to be processed in real-time, and is simple enough to be searched and measured on.
As a result, an Original Frame is created, and a new column extracted from the Original Frame is added to the Original Timeline.

\begin{figure}[!h]
	\centering
	\scalebox{0.5}{\input{figures/frame_transforms.pdf_tex}}
	\caption{Preprocessing steps: the series of transformations performed on each frame before detection. The goal is to create an effective frame format for the following processing steps.\label{fig:transforms}}
\end{figure}

The first stage is perspective compensation.
At this point, the image, distorted by perspective transform, is remapped based on a calibration rectangle, so that lengths on the frame become independent from the distance of the camera.
The later performed vehicle-size calculation and classification strongly relies on the assumption that vehicle-lengths are not subject to perspective distortion, and are irrelative to their position on the frame.
The calibration rectangle is defined manually using the Project Configurator user interface (UI)\ref{subs:ProjectConfigurator}.

The second stage is rotation of the frame, so that the tripwire becomes vertical.
This form allows measurements in the dimension perpendicular to the tripwire to be performed by reading solely certain rows of the image.
This method, considering the image-storage technique of the OpenCV library, where frames are stored as an array of rows in the memory, speeds up the computation significantly, and simplifies searching on frames.

In the third level of preprocessing a resizing is performed in order to decrease the number of pixels and the computation time, and increase the speed of processing.
The end-size of the frame can be specified through the configuration files, and is usually 320 by 240 pixels.
%----------------------------------------------------------------------------
\subsection{Background-subtraction}
%----------------------------------------------------------------------------
The second phase of processing is background-subtraction, that is the separation of moving objects form the static background on each frame.
This step creates the MOG2 Frame and Timeline as a result of background-subtraction.
The correction of the subtraction method use the Motionless Timeline.

This step is performed by a GMM-based method, Mixture of Gaussians 2 (MOG2), that is available in the OpenCV library.
This method associates each pixel with a mixture of Gaussian intensity-distributions, and identifies each new data as background, if it is part of the constructed distribution-model, or as a moving object, if it falls out of its range.

The drawback of this technique is its complexity, that comes with an increased computational cost and time.
Throughout the processing, background subtraction has the highest computation time, causing it to be the most computationally expensive of all of the outlined processing steps.
However, since the following steps rely on the accuracy of the result of the background-subtraction, the robustness and precision of the approach are essential.

Another drawback of background-modelling is that temporarily stopped objects that have been motionless for a while are missed, as they are being included into the background model.
The errors caused by this is corrected later on, by detecting immobile vehicles.

Subtraction is performed on the whole frame, although some approaches model only the tripwire's pixels as background--foreground.
That is, because features, such as size and velocity are calculated by measuring the foreground-object's lengths that are perpendicular to the tripwire on the MOG2 Frame.

After the background-subtraction post-corrections are performed, in order to filter some errors on the MOG2 output.

First, morphological opening followed by closing removes punctual noises, and smooths object contours.
As a result, the Morphology Timeline is created.
After, stopped vehicles are detected, using the Motionless Timeline.
As sudden light changes can cause an overreaction of MOG2, to prevent this, the MOG2 Frames are filtered.
If a the count of a frame's white pixels exceeds a limit, the frame is cleared, and switched to background, to avoid false detections.
Although this process causes some vehicles to be missed, empirical analysis shows, that error rate is reasonably lower if the overreaction frames are filtered.
%----------------------------------------------------------------------------
\subsection{Data extraction}
%----------------------------------------------------------------------------
After the MOG2 Timeline and Frame are created, and pixels of moving objects are identified, data extraction is performed.
First, vehicles are detected with connected component analysis on the MOG2 Timeline image.
Since some vehicles can fuse on the timeline image, if they reach the tripwire together, the result of counting is
%----------------------------------------------------------------------------
\subsection{Evaluation}
%----------------------------------------------------------------------------
\subsubsection{The parameter evaluation process}
%----------------------------------------------------------------------------





%----------------------------------------------------------------------------
\chapter{Hardware and software environment}\label{chap:Environment}
%----------------------------------------------------------------------------
The primary task of the SOLSUN Traffic Sensor is counting, that is performed by the core software.
- hw and supplementary sw elements facilitate the operation of the sensor in the production environment
- ssw is responsible for connecting the sw to the surrounding environment transfers the user input and output
%----------------------------------------------------------------------------
\section{Supplementary software elements}\label{sec:SupplementarySoftware}
%---------------------------------------------------------------------------
- user input: settings
- facilitate functioning under heterogeneous circumstances
- and the testing
- evaluation and diagnostics
- provide data to the user, generates user output: 
- visualizes the data stored in the system in the form of timeline images, fram-strips and other container objects 
- results of system diagnostics, 
- the the result of counting and classification in different representations
- for human understanding
- these element help to test and evaluate the core software
- and connect to the EnTalk Network
%----------------------------------------------------------------------------
\subsection{Configuration process}\label{subs:ProjectConfigurator}
%----------------------------------------------------------------------------
- various algorithm parameters
- possibility of switch proessing steps

\subsubsection{Calibration process}
%----------------------------------------------------------------------------
Manually thorugh the project Configurator
- tripwire
- calibration rectangle
- optional road lanes
- the beginning and ending frame index of the video to be processed
%----------------------------------------------------------------------------
\subsection{Evaluation}
%----------------------------------------------------------------------------

\subsubsection{Ground truth}


%----------------------------------------------------------------------------
\subsection{Data representation}
%----------------------------------------------------------------------------
\subsubsection{Timeline Images}

\subsubsection{Videos}
%----------------------------------------------------------------------------
\subsection{Operation modes of the Traffic Sensor}
%----------------------------------------------------------------------------
\subsubsection{Single-run or test mode} 
- from video file or camera stream with a specified amount of frames 
- at the end of the operation, evaluation is possible, 
- entire timeline images can be saved, 
- demo video creation is possible and the recording of the raw video stream is possible

\begin{figure}[!h]
	\centering
	\includesvg[width=\textwidth]{process_run_types}
	\caption{. \label{fig:run_types}}
\end{figure}

\subsubsection{Continuous run}
- operation with a video stream from a camera (works with files as well for testing purposes)
- timeline images can be saved only in pieces
- evaluation and other features that is not possible, because there is no definite end to the process
%----------------------------------------------------------------------------
\section{Hardware environment}
%----------------------------------------------------------------------------

\subsection{The board}

\subsection{Network architecture}


%----------------------------------------------------------------------------


%----------------------------------------------------------------------------
\chapter{Experimental results}\label{chap:Tests}
%----------------------------------------------------------------------------
In the following chapter the results of the system tests will be discussed. 
Performance and precision criteria 
%----------------------------------------------------------------------------
\section{Evaluation results}
%----------------------------------------------------------------------------
\subsection{Test videos}


\subsection{Results}
\subsubsection{Initial setting}
- minSize: 250
- lanes: defined for most videos

\subsubsection{The minimal size of blobs redefined}
- minSize: 200
- lanes: defined for most videos

\subsubsection{No lanes defined}
- minSize: 200
- lanes: no lanes

\subsubsection{Split size added}
- minSize: 100
- splitMinSize : 100

\subsection{Common errors and solutions}
\subsubsection{False positives}
\subsubsection{False negatives}
\subsubsection{Misclassification errors}

%----------------------------------------------------------------------------
\section{Calibration correction}\label{chap:cal_corr}
%----------------------------------------------------------------------------
Due to perspective distortion, vehicles, that in effect have the same sizes differ in length on the Raw Frames.
Since vehicle lengths are the basis of classification, perspective distortion is eliminated using a quadrangle, called the calibration rectangle, that is defined by the user.

The calibration rectangle is hand-drawn on a frame in the Project Configurator GUI (section \ref{chap:calibration}.) as part of the calibration process.
The position and the sides' lengths are estimated with the help of any occurring road lanes, or simply based on naked eye examination.
On this wise, due to the inaccuracy of human perception, the location and size of the calibration rectangle are significantly uncertain.
Empirical analysis shows that errors are actually present in many hand-calibrated videos.

To avoid manual calibration, some traffic surveillance systems automatize this process.
However in heterogeneous environment, it is a highly complex problem, and is not implemented in the current version of the Traffic Sensor system.

To correct the existent settings, the extent of the miscalibration is evaluated examining the length of the vehicles close to and far from the camera.
This way the calibration rectangle is redefined for later use during the evaluation phase of the processing.

\subsection{The miscalibration problem}
\subsubsection{Perspective compensation method}
The calibration rectangle is constructed by defining its four edges from the upper left corner, to the lower left one in order.
The edges must correspond to a square's in the 3D world, that due to the perspective transformation, show as a general rectangle on the image plane.
On most frames, lane markings are used for this purpose.

The transform projects the four vertices back into a square, along with the entire pixel grid of the image using a transformation map to speed up the process.
For defining the transform and constructing the map two universal OpenCV methods were used\cite{PersTrans, WrapPers}.

In the result images -- the Original Frame -- the vehicle lengths are independent of the position in the image plane. 
This opens the possibility of correct calibration.

The calibration rectangle's sides are also used to define a scaling factor between the 3d world and the image plane.
This can convert all lengths in the image to a metric scale.

\subsubsection{Miscalibrated streams}
- if the edges of the rectangle do not fit a perfect square in reality (valosag)
- the vehicles on the compensated frames are not the same length
- in most cases their lengths still decrease as (távolodva) from the camera
- easily (felismerheto) (jelenseg): the int the farest?? road lane cares are classified as bicycles

\subsection{Solutions}
\begin{enumerate}
	\item  \textbf{Estimate decrease}
	\item  \textbf{Calculate proper size ratios}
	\item  \textbf{Redraw calibration rectangle}
\end{enumerate}
- quantify the decrease of the lengths
- based on the correct lengths, redraw the calibration rect

\subsection{Calibration correction results}

%----------------------------------------------------------------------------
\section{Speed testing}
%----------------------------------------------------------------------------
- Speed criteria to meet:
	- minimum speed \SI{25}{frame/s}
	- loosing frames and thus vehicles
	- the optimum is higher, \SI{30}{frame/s}
	- important to be higher for the integration of new functions, processing steps, or extending existing ones
	- e.g. shadow-detection feature, or a more sophisticated classifier

\subsection{Results}

\renewcommand{\thechapter}{\Alph{chapter}}
\setcounter{chapter}{3}  % a fofejezet-szamlalo az angol ABC 6. betuje (F) lesz
\chapter*{Conclusions}\addcontentsline{toc}{chapter}{Conclusions}\label{chap:Conclusions}
\setcounter{equation}{0} 
\setcounter{section}{0}
\numberwithin{equation}{section}
\numberwithin{figure}{section}
\numberwithin{lstlisting}{section}
\section{Summary}
In this work a software system created for counting and classifying vehicles in urban scenes using image processing algorithms is presented.

The software was implemented as an end product of the Traffic Sensor sub-project, that is part of the larger SOLSUN project.
The aim of the main project is to devise an intelligent sensor system, built into a street-lighting lamp, that is able to monitor the environment with various sensors at the same time.
The data collected by the sensors are used to estimate pollution and reduce greenhouse gas emission and energy consumption. 
The Traffic Sensor sub-project develops the sensor that is responsible for the surveillance of traffic by recording the road with a camera and estimating the count of vehicles passing by.
The project's requirements include a pre-defined embedded target hardware with passive cooling ability, real-time operation of the software, and a counting precision of at least \SI{70}{\%}.

The work outlines the overall structure and operation procedure of the Traffic Sensor system. 

It introduces an effective image domain reduction paradigm, the creation of Timeline Images, that is the basic idea of the processing.
The thesis describes the generation procedure of the Timeline Images, and their role in the operation of the Sensor.
Besides that, the core software, that is responsible for the main functionality, is presented.
The architectural structure and operation of the software is described in detail, including the video processing steps, that consist of the transform of frames, the detection of moving objects with background subtraction, calculation of traffic parameters, and classification of the detected vehicles using their length.
For each processing step the used TIs and their role is described as well.

The environment of the core system is also summarised.
The target hardware and the network architecture responsible for data transmission, and the software parts that connect the Traffic Sensor to its surroundings are described.
Other supplementary functionalities created for system diagnostic purposes and performance evaluation are detailed as well.

The system was tested with regard to its precision and real-time ability characteristics to fit the pre-defined criteria.
The test results are presented and options are proposed and analysed to improve the precision.
An algorithm for the correction and detection of inaccurate calibration is introduced, and tested.

\section{My own contribution within the context of the whole project}
As it was mentioned in the introduction (\ref{sec:contribution_of_the_thesis}), this work presents the structure of the entire Traffic Sensor system as a whole, although it is the product of collective work, and not exclusively my own.
However my own contributions are emphasized and presented in more detail throughout the thesis.
For clarification purposes, in this section, contributors of each part of the project are briefly summarised.

\noindent Parts of the system that were the result of the team's collective work are as follows:
\begin{itemize}
	\item[-] the overall software architecture of the system (section \ref{sec:software_architecture})
	\item[-] the concrete steps of the operation procedure (section \ref{sec:processing_steps}).
\end{itemize}

\noindent Results that are the product of my own exclusive work are as follows:
\begin{itemize}
	\item[-] the implementation of some of the Processors, including the ones responsible for background-subtraction and morphological correction (section \ref{sec:background-subtraction})
	\item[-] the development and implementation of the parameter extraction algorithm (section \ref{chap:parameter_evaluation})
	\item[-] the development and implementation of the classifier (section \ref{sec:classification})
	\item[-] realization of the visualization and diagnostics features (section \ref{sec:system_diagnostics})
	\item[-] the implementation of the evaluation process and devising the performance measures (section \ref{chap:evaluation})
	\item[-] the analysis of the system's performance (section \ref{sec:evaluation_results})
	\item[-] the development and implementation of the calibration correction algorithm (section \ref{sec:calibration_correction_method}).
\end{itemize}
For further details, see my other paper on the topic \cite{Barancsuk2016}.

\noindent Other system parts only mentioned in the thesis for a coherent presentation, that were developed by other members of the team are as follows:
\begin{itemize}
	\item[-] the procedure of porting the software to the target hardware
	\item[-] the software parts responsible for connecting with the environment and configuration features (section \ref{sec:communication})
	\item[-] the continuous integration of the system (section \ref{sec:continuous_integration})
	\item[-] the recording of the test videos (section \ref{sec:test_videos})
	\item[-] the development and implementation of the blob splitting and fusing algorithms (section \ref{sec:background-subtraction})
	\item[-] the development and implementation of the preprocessing phase (section \ref{sec:preprocessing}).
\end{itemize}

%The tasks detailed in the specification of the thesis work were complied.

\section{Perspectives}
Although the system complies with the criteria presented above, the performance measures can be further improved by filtering and correcting common errors that are still present in the detection results.

One of the main error causes is the presence of shadows, that, although are partly removed by the built-in method of the MOG2 background-subtractor, can remain unfiltered and are the still the primary sources of false detections in shadow-heavy videos.
Many shadow-removal algorithms are introduced in literature, from which a suitable one could be integrated into the processing.

Counting at night is considerably reliable at the current version of the Traffic Sensor, although it can be further improved by developing an algorithm for identifying and removing blobs caused by the headlights of cars.

Another progression direction could be the improvement of the existing blob splitting and fusing methods, that in the current version tend to fail with uncommon cases and in complex surroundings.

To advance the classification result, a more sophisticated classifier, e.g. one deciding based on a parameter vector can be implemented.

Another weak point of the system is that it does not handle extreme, exceptional weather conditions, like fog, rain or snow.
Detection in these circumstances is a highly complex problem, and is still a heavily researched topic, however these weather conditions can be easily identified, and the Sensor could be temporarily turned off until reliable detection is possible again.
	
As a conclusion, the performance of the counting could only be improved by integrating more sophisticated algorithms into the processing steps, however the complexity of these methods would make real-time operation impossible on the current target hardware.
%----------------------------------------------------------------------------
\chapter*{Acknowledgement}\addcontentsline{toc}{chapter}{Acknowledgement}
%----------------------------------------------------------------------------

I would like to offer my special thanks to my advisor, Dr.~Kristóf Csorba for
I would like to express my gratitude to Dr.~László Blázovics

I would like to express my very great appreciation to ***
I would like to offer my special thanks to ***
Advice given by *** has been a great help in ***
I am particularly grateful for the assistance given by ***
Assistance provided by *** was greatly appreciated.
I wish to acknowledge the help provided by ***
Dr Bilicz provided me with very valuable ***.
I would like to thank the following companies for their assistance with the collection of my data:
***
***
My special thanks are extended to the staff of *** company for ***

\begin{displaymath}
	\star \star \star
\end{displaymath}
%----------------------------------------------------------------------------
\appendix
%----------------------------------------------------------------------------
\chapter*{Appendices}\addcontentsline{toc}{chapter}{Appendices}\label{chap:Appendices}
\setcounter{chapter}{1}  % a fofejezet-szamlalo az angol ABC 6. betuje (F) lesz
\setcounter{equation}{0} 
\numberwithin{equation}{section}
\numberwithin{figure}{section}
\numberwithin{lstlisting}{section}
%\numberwithin{tabular}{section}

%----------------------------------------------------------------------------
\section{Advertisement application}
%----------------------------------------------------------------------------
The advertisement application is an alternative utilization of the Traffic Sensor's functionality, availing of that the Sensor is able to detect and count any moving objects in its environment.

The goal of the application is to estimate the number of people present on the corridor of the \viktanszek, where the Sensor will be deployed.
The Traffic Sensor system will be connected to a screen with a web interface, that shows miscellaneous advertisements.

Advertisers will be able to buy time-intervals for their advertisements for different prices.
The intervals' prices are assessed based on the count of people present at that time.

The Traffic Sensor is already fit for the task, however the following settings were modified for the changed situation.
\begin{enumerate}
	\item The camera angle is horizontal, thus the Tripwire is set to vertical.
	
	\item Since the bypassing people are close to the camera, they are not subject to significant  perspective distortion, and there is no need for perspective compensation.
	
	\item The other consequence of the closeness of the camera is that moving objects fill most of the frame, and to detect them, only a rather small Tripwire is necessary, approximately a \SI{20}{pixel} long section.
	An advantage resulting from this is the low computation time.
	
	\item The minimal size of blobs is adjusted to the size of the detected objects empirically.
	
	\item There is no need for classification, since it is highly likely, that on the corridor only people are present.
\end{enumerate}

Although the Advertisement Application is still in pilot phase, and is not fully deployed yet, some example images of the configuration process are seen in figures \ref{fig:adv_app_config}., \ref{fig:adv_app_full_orig}., \ref{fig:adv_app_counting}.

\begin{figure}[!h]
	\centering
	\includegraphics[width=0.5\textwidth]{ad_app_config.png}
	
	\caption{The GUI of the Project Configurator, showing the initial configuration of the Advertisement Application. Here, the Tripwire crosses the whole frame, although later it is decreased, since the application can count people using a much smaller portion of the line.  \label{fig:adv_app_config}}
\end{figure}

\begin{figure}[!h]
	\centering
	\includegraphics[width=0.6\textwidth]{orig_full.png}
	
	\caption{Part of the Original Timeline Image showing a series of people walking in front of the camera. Since the person occupies the whole height of the image, for detection only a smaller section of the image is enough, as seen on figure \ref{fig:adv_app_counting}.\label{fig:adv_app_full_orig}}
\end{figure}

\begin{figure}[!h]
	\centering
	\begin{subfigure}[t]{0.6\textwidth}
	\includegraphics[width=\textwidth]{original_section.png}
	\subcaption{Part of the Original Timeline Image.}
	\end{subfigure}
	\hfill
	\begin{subfigure}[t]{0.6\textwidth}
		\includegraphics[width=\textwidth]{mog2_section.png}
		\subcaption{Part of the MOG2 Timeline Image.}
	\end{subfigure}

	\caption{Timeline Images of the Advertisement Application. The MOG2 image captures the moving object correctly each time, but also some noise due to light changes.\label{fig:adv_app_counting}}
\end{figure}

%----------------------------------------------------------------------------
\clearpage\section{VideoMixer application}
%----------------------------------------------------------------------------
The VideoMixer application was created as a side-project for the customers of the Traffic Sensor.

The software is responsible for the visualization of multiple video streams at the same time for demonstration purposes.

The videos can either be recorded from a camera or read from a file.
The user can control, which videos are seen at a certain time, that are blended by the application on top of each other.
The speed of the playing and the looping of the videos can be set as well.

\begin{figure}[!h]
	\centering
	\includegraphics[width=0.5\textwidth]{vm_menu.png}
	
	\caption{The menu of the VideoMixer. Whenever a key is pressed, it is displayed, as well as the possible control options. \label{fig:video_mixer_menu}}
\end{figure}

The settings (e.g the name and type of the input videos, the controlling keys and the whether the videos are looped) are defined in a configuration file, similarly to the Traffic Sensor's (seen in code-listing \ref{lst:config_vm}).

The application can be utilized as a demonstration of the Traffic Sensor's operation as well, some example frames are figures \ref{fig:video_mixer_frames}., \ref{fig:video_mixer_menu}..

\begin{lstlisting}[frame=single,float=!ht,caption={Part of a configuration file for the VideoMixer application. The file sets parameters for the main input video (section [mainVideoInput]), and the overlay videos (section [overlayVideos]), that are blended on top. The file defines whether the video is read from a camera or a file, whether some parts of the background are removed, and other adjustments, like the comment displayed on the screen, when a video starts playing.},label=lst:config_vm]
{
[mainVideoInput]
cameraInput = false
cameraIndex = 0
filename = orig.avi
loop = true
removeBackgroundColor = 0x00ff00

[overlayVideos]
filename_0 = mog2.avi
gpioNum_0 = 0
overlayRatio_0 = 0.4
comment_0 = playing file 1
displayCommentOnScreen_0 = false
removeBackgroundColor_0 =  0x00ff00
loop_0 = true
filename_1 = speed.avi
...
\end{lstlisting}

\begin{figure}[!h]
	\centering
	\begin{subfigure}[t]{0.6\textwidth}
		\includegraphics[width=\textwidth]{orig_mog22.png}
		\subcaption{The overlay video captures MOG2 Frames and the building process of a MOG2 Timeline Image.}
	\end{subfigure}
	\hfill
	\begin{subfigure}[t]{0.6\textwidth}
		\includegraphics[width=\textwidth]{orig_result.png}
		\subcaption{The overlay video captures a Result Timeline image, that shows the detected vehicles with different colors projected on top of the main video. The Result Timeline also displays vehicle parameters, such as size and type. Blending this videos together, both correct and faulty detections can be analysed.}
	\end{subfigure}
	\hfill
	\begin{subfigure}[t]{0.6\textwidth}
		\includegraphics[width=\textwidth]{orig_speeds.png}
		\subcaption{The blended video is a Speed Timeline. The direction of each detected vehicle can be read by the colors of the Speed Timeline (with red the vehicles moving from the right to the left are marked, and the ones heading in the other direction are colored with blue). }
	\end{subfigure}

	\caption{Frames of the VideoMixer application. The overlay video is blended on top of the main video, that consists of an Original Frame and an Original Timeline recorded using the Traffic Sensor. Overlay videos are other Timeline Images, MOG2, Speed and Result.\label{fig:video_mixer_frames}}
\end{figure}

%----------------------------------------------------------------------------
\clearpage\section{Some examples of the output Timeline Images}\label{app:TIs}
%----------------------------------------------------------------------------
This appendix presents some test results in the form of output Timeline Images, focusing on remarkable points and interesting details.

\centering
\begin{figure}[h!]
	\centering
	\captionsetup{width=0.9\textwidth,aboveskip=10pt}
	\includegraphics[width=\textwidth]{GoldA1.png}
	\caption[Some output Timelines of the GoldA1 test video]{The figures (from the top to the bottom) are as follows: Original Timeline Image, MOG2 Timeline Image, Size Timeline Image, Speed Timeline Image and Evaluation Result Timeline Image respectively. \\
	On night-time videos like GoldA1 a prevalent error cause is that headlights of vehicles, that are reflected back by the road lead to false alarms (some examples are on the right side of the image). Headlight blobs fusing with the following vehicles are also sources of misclassification errors, since the vehicle size increases (seen on the Size Timeline). \\
	One can notice that the colours on the Speed Timeline  mark the directions of the road lanes. \label{fig:GoldA1}}
\end{figure}

\centering
\begin{figure}[h!]
	\centering
	\captionsetup{width=0.9\textwidth,aboveskip=10pt}
	\includegraphics[width=0.8\textwidth]{GoldB12.png}
	\caption[Some output Timelines of the GoldB1 test video]{The figures (from the top to the bottom) are as follows: Original Timeline Image, MOG2 Timeline Image, Motionless Timeline Image, Size Timeline Image, Speed Timeline Image and Evaluation Result Timeline Image respectively.\\
	 The following common error cases are seen in the Evaluation Result Timeline Image. Small blobs of motorcycles are filtered out (in the middle crossed with red) generating false negative detections. Some stopping vehicles add noise to the detection when missed by the fusing algorithm (false negatives caused by noise are on the left side of the image). The splitting mechanism tend to erroneously cut vehicles, that have a peculiar shape (in the middle and in the right). \\
	 One can notice, that the Parameter Timelines were calculated without lanes defined, since the whole outlines of blobs are filled with values. Intensities on the Size Timeline shows well the approximate length of vehicles. \label{fig:GoldB1}}
\end{figure}

\clearpage
\centering
\begin{sidewaysfigure}[h!]
	\centering
	\captionsetup{width=0.9\textwidth,aboveskip=10pt}
	\includegraphics[width=\textwidth]{ErE5.png}
	\caption[Some output Timelines of the ErE5 test video]{The figures (from the top to the bottom) are as follows: Original Timeline Image, MOG2 Timeline Image, Size Timeline Image, Following Distance Timeline Image, Speed Timeline Image and Result Timeline Image respectively.\\
    The video was recorded at the IBikeBudapest cyclist demonstration, therefore road users are almost exclusively cyclists. \\
    Comparing the Size and the Following Distance Timeline, it is observable, that following distance and length are calculated in alternate times, whether the Tripwire is empty (following distances) or occupied by a vehicle (length). An interesting point of the result is cars appearing in red on the top part of the Speed Timeline, heading in the opposite direction as the cyclists. Different coloured dots inside the vehicles in the Speed Timeline are caused by noise, that is the result of the size changes of the vehicle (further details on parameter calculation based on vehicle position is in chapter \ref{chap:parameter_evaluation}.). \label{fig:ErE5}}
\end{sidewaysfigure}

\centering
\begin{sidewaysfigure}[h!]
	\centering
	\captionsetup{width=0.9\textwidth,aboveskip=10pt}
	\includegraphics[width=\textwidth]{Fovam_C.png}
	\caption[Some output Timelines of the Fovam\_C test video]{The figures (from the top to the bottom) are as follows: Original Timeline Image, MOG2 Timeline Image and Result Timeline Image respectively. \\
	The video was recorded at sundown, and captures only a few cars. Although the environmental light frequently changes (an example is the lilac section on the left on the Original Timeline), the count of false detections is low. Fused cars in the middle are split appropriately. However the stopping pedestrian (top right corner) is ignored, al well as the other two pedestrians who are falsely split based on their shapes.}
\end{sidewaysfigure}

\centering
\begin{sidewaysfigure}[h!]
	\centering
	\captionsetup{width=0.9\textwidth,aboveskip=10pt}
	\includegraphics[width=\textwidth]{ErC1.png}
	\caption[Some output Timelines of the ErC1 test video]{The figures (from the top to the bottom) are as follows: Original Timeline Image, Result Timeline Image and Evaluation Result Timeline Image respectively. The video captures a stopping tram, a handful of cars and pedestrians. The tram's slow movement is still a challenge to the fusing algorithm, thus causes some noise.}%{ErC kódú videó futtatási eredménye. Fentről lefelé a képek rendre a következők. entről lefelé a képek rendre a következők. Eredeti videóból képzett TK. MOG2 TK. A kiértékelés eredménye. \\
		%A videón egy megálló villamos, illetve a mellette haladó gyalogosok és autók láthatók. A rendszer jól összeilleszti a megálló villamos blobját, de helytelenül a jobb alsó sarokban található blobokat (hosszú téglalap), amely csak a a háttér folytonossága miatt tűnik a rendszernek megálló autónak.}
\end{sidewaysfigure}

\centering
\begin{sidewaysfigure}[h!]
	\centering
	\captionsetup{width=0.9\textwidth,aboveskip=10pt}
	\includegraphics[width=\textwidth]{SasA1.png}
	\caption[Some output Timelines of the SasA1 test video]{The figures (from the top to the bottom) are as follows: Original Timeline Image, MOG2 Timeline Image, Speed Timeline Image and Evaluation Result Timeline Image respectively. \\
	The software performs well on this video. Most fused vehicles are correctly split, and stopping vehicles' blobs are reattached as well. Due to shaky camera, that causes the background to change, some false alarms are present (on the top, middle part). The Speed Timeline is calculated only in the centre of road lanes displaying fewer values than the Parameter Timelines in figures \ref{fig:ErE5}., \ref{fig:GoldA1}., \ref{fig:GoldB1}.}
\end{sidewaysfigure}




\newglossaryentry{pqr}{name=PQR, description={Coordinated Universal Time}}
%\listoftables\addcontentsline{toc}{chapter}{List of tables}

\bibliography{trafficVideoSens}
\addcontentsline{toc}{chapter}{Bibliography}
\bibliographystyle{ieeetr}

%%----------------------------------------------------------------------------
\appendix
%----------------------------------------------------------------------------
\chapter*{Appendices}\addcontentsline{toc}{chapter}{Appendices}\label{chap:Appendices}
\setcounter{chapter}{1}  % a fofejezet-szamlalo az angol ABC 6. betuje (F) lesz
\setcounter{equation}{0} 
\numberwithin{equation}{section}
\numberwithin{figure}{section}
\numberwithin{lstlisting}{section}
%\numberwithin{tabular}{section}

%----------------------------------------------------------------------------
\section{Advertisement application}
%----------------------------------------------------------------------------
The advertisement application is an alternative utilization of the Traffic Sensor's functionality, availing of that the Sensor is able to detect and count any moving objects in its environment.

The goal of the application is to estimate the number of people present on the corridor of the \viktanszek, where the Sensor will be deployed.
The Traffic Sensor system will be connected to a screen with a web interface, that shows miscellaneous advertisements.

Advertisers will be able to buy time-intervals for their advertisements for different prices.
The intervals' prices are assessed based on the count of people present at that time.

The Traffic Sensor is already fit for the task, however the following settings were modified for the changed situation.
\begin{enumerate}
	\item The camera angle is horizontal, thus the Tripwire is set to vertical.
	
	\item Since the bypassing people are close to the camera, they are not subject to significant  perspective distortion, and there is no need for perspective compensation.
	
	\item The other consequence of the closeness of the camera is that moving objects fill most of the frame, and to detect them, only a rather small Tripwire is necessary, approximately a \SI{20}{pixel} long section.
	An advantage resulting from this is the low computation time.
	
	\item The minimal size of blobs is adjusted to the size of the detected objects empirically.
	
	\item There is no need for classification, since it is highly likely, that on the corridor only people are present.
\end{enumerate}

Although the Advertisement Application is still in pilot phase, and is not fully deployed yet, some example images of the configuration process are seen in figures \ref{fig:adv_app_config}., \ref{fig:adv_app_full_orig}., \ref{fig:adv_app_counting}.

\begin{figure}[!h]
	\centering
	\includegraphics[width=0.5\textwidth]{ad_app_config.png}
	
	\caption{The GUI of the Project Configurator, showing the initial configuration of the Advertisement Application. Here, the Tripwire crosses the whole frame, although later it is decreased, since the application can count people using a much smaller portion of the line.  \label{fig:adv_app_config}}
\end{figure}

\begin{figure}[!h]
	\centering
	\includegraphics[width=0.6\textwidth]{orig_full.png}
	
	\caption{Part of the Original Timeline Image showing a series of people walking in front of the camera. Since the person occupies the whole height of the image, for detection only a smaller section of the image is enough, as seen on figure \ref{fig:adv_app_counting}.\label{fig:adv_app_full_orig}}
\end{figure}

\begin{figure}[!h]
	\centering
	\begin{subfigure}[t]{0.6\textwidth}
	\includegraphics[width=\textwidth]{original_section.png}
	\subcaption{Part of the Original Timeline Image.}
	\end{subfigure}
	\hfill
	\begin{subfigure}[t]{0.6\textwidth}
		\includegraphics[width=\textwidth]{mog2_section.png}
		\subcaption{Part of the MOG2 Timeline Image.}
	\end{subfigure}

	\caption{Timeline Images of the Advertisement Application. The MOG2 image captures the moving object correctly each time, but also some noise due to light changes.\label{fig:adv_app_counting}}
\end{figure}

%----------------------------------------------------------------------------
\clearpage\section{VideoMixer application}
%----------------------------------------------------------------------------
The VideoMixer application was created as a side-project for the customers of the Traffic Sensor.

The software is responsible for the visualization of multiple video streams at the same time for demonstration purposes.

The videos can either be recorded from a camera or read from a file.
The user can control, which videos are seen at a certain time, that are blended by the application on top of each other.
The speed of the playing and the looping of the videos can be set as well.

\begin{figure}[!h]
	\centering
	\includegraphics[width=0.5\textwidth]{vm_menu.png}
	
	\caption{The menu of the VideoMixer. Whenever a key is pressed, it is displayed, as well as the possible control options. \label{fig:video_mixer_menu}}
\end{figure}

The settings (e.g the name and type of the input videos, the controlling keys and the whether the videos are looped) are defined in a configuration file, similarly to the Traffic Sensor's (seen in code-listing \ref{lst:config_vm}).

The application can be utilized as a demonstration of the Traffic Sensor's operation as well, some example frames are figures \ref{fig:video_mixer_frames}., \ref{fig:video_mixer_menu}..

\begin{lstlisting}[frame=single,float=!ht,caption={Part of a configuration file for the VideoMixer application. The file sets parameters for the main input video (section [mainVideoInput]), and the overlay videos (section [overlayVideos]), that are blended on top. The file defines whether the video is read from a camera or a file, whether some parts of the background are removed, and other adjustments, like the comment displayed on the screen, when a video starts playing.},label=lst:config_vm]
{
[mainVideoInput]
cameraInput = false
cameraIndex = 0
filename = orig.avi
loop = true
removeBackgroundColor = 0x00ff00

[overlayVideos]
filename_0 = mog2.avi
gpioNum_0 = 0
overlayRatio_0 = 0.4
comment_0 = playing file 1
displayCommentOnScreen_0 = false
removeBackgroundColor_0 =  0x00ff00
loop_0 = true
filename_1 = speed.avi
...
\end{lstlisting}

\begin{figure}[!h]
	\centering
	\begin{subfigure}[t]{0.6\textwidth}
		\includegraphics[width=\textwidth]{orig_mog22.png}
		\subcaption{The overlay video captures MOG2 Frames and the building process of a MOG2 Timeline Image.}
	\end{subfigure}
	\hfill
	\begin{subfigure}[t]{0.6\textwidth}
		\includegraphics[width=\textwidth]{orig_result.png}
		\subcaption{The overlay video captures a Result Timeline image, that shows the detected vehicles with different colors projected on top of the main video. The Result Timeline also displays vehicle parameters, such as size and type. Blending this videos together, both correct and faulty detections can be analysed.}
	\end{subfigure}
	\hfill
	\begin{subfigure}[t]{0.6\textwidth}
		\includegraphics[width=\textwidth]{orig_speeds.png}
		\subcaption{The blended video is a Speed Timeline. The direction of each detected vehicle can be read by the colors of the Speed Timeline (with red the vehicles moving from the right to the left are marked, and the ones heading in the other direction are colored with blue). }
	\end{subfigure}

	\caption{Frames of the VideoMixer application. The overlay video is blended on top of the main video, that consists of an Original Frame and an Original Timeline recorded using the Traffic Sensor. Overlay videos are other Timeline Images, MOG2, Speed and Result.\label{fig:video_mixer_frames}}
\end{figure}

%----------------------------------------------------------------------------
\clearpage\section{Some examples of the output Timeline Images}\label{app:TIs}
%----------------------------------------------------------------------------
This appendix presents some test results in the form of output Timeline Images, focusing on remarkable points and interesting details.

\centering
\begin{figure}[h!]
	\centering
	\captionsetup{width=0.9\textwidth,aboveskip=10pt}
	\includegraphics[width=\textwidth]{GoldA1.png}
	\caption[Some output Timelines of the GoldA1 test video]{The figures (from the top to the bottom) are as follows: Original Timeline Image, MOG2 Timeline Image, Size Timeline Image, Speed Timeline Image and Evaluation Result Timeline Image respectively. \\
	On night-time videos like GoldA1 a prevalent error cause is that headlights of vehicles, that are reflected back by the road lead to false alarms (some examples are on the right side of the image). Headlight blobs fusing with the following vehicles are also sources of misclassification errors, since the vehicle size increases (seen on the Size Timeline). \\
	One can notice that the colours on the Speed Timeline  mark the directions of the road lanes. \label{fig:GoldA1}}
\end{figure}

\centering
\begin{figure}[h!]
	\centering
	\captionsetup{width=0.9\textwidth,aboveskip=10pt}
	\includegraphics[width=0.8\textwidth]{GoldB12.png}
	\caption[Some output Timelines of the GoldB1 test video]{The figures (from the top to the bottom) are as follows: Original Timeline Image, MOG2 Timeline Image, Motionless Timeline Image, Size Timeline Image, Speed Timeline Image and Evaluation Result Timeline Image respectively.\\
	 The following common error cases are seen in the Evaluation Result Timeline Image. Small blobs of motorcycles are filtered out (in the middle crossed with red) generating false negative detections. Some stopping vehicles add noise to the detection when missed by the fusing algorithm (false negatives caused by noise are on the left side of the image). The splitting mechanism tend to erroneously cut vehicles, that have a peculiar shape (in the middle and in the right). \\
	 One can notice, that the Parameter Timelines were calculated without lanes defined, since the whole outlines of blobs are filled with values. Intensities on the Size Timeline shows well the approximate length of vehicles. \label{fig:GoldB1}}
\end{figure}

\clearpage
\centering
\begin{sidewaysfigure}[h!]
	\centering
	\captionsetup{width=0.9\textwidth,aboveskip=10pt}
	\includegraphics[width=\textwidth]{ErE5.png}
	\caption[Some output Timelines of the ErE5 test video]{The figures (from the top to the bottom) are as follows: Original Timeline Image, MOG2 Timeline Image, Size Timeline Image, Following Distance Timeline Image, Speed Timeline Image and Result Timeline Image respectively.\\
    The video was recorded at the IBikeBudapest cyclist demonstration, therefore road users are almost exclusively cyclists. \\
    Comparing the Size and the Following Distance Timeline, it is observable, that following distance and length are calculated in alternate times, whether the Tripwire is empty (following distances) or occupied by a vehicle (length). An interesting point of the result is cars appearing in red on the top part of the Speed Timeline, heading in the opposite direction as the cyclists. Different coloured dots inside the vehicles in the Speed Timeline are caused by noise, that is the result of the size changes of the vehicle (further details on parameter calculation based on vehicle position is in chapter \ref{chap:parameter_evaluation}.). \label{fig:ErE5}}
\end{sidewaysfigure}

\centering
\begin{sidewaysfigure}[h!]
	\centering
	\captionsetup{width=0.9\textwidth,aboveskip=10pt}
	\includegraphics[width=\textwidth]{Fovam_C.png}
	\caption[Some output Timelines of the Fovam\_C test video]{The figures (from the top to the bottom) are as follows: Original Timeline Image, MOG2 Timeline Image and Result Timeline Image respectively. \\
	The video was recorded at sundown, and captures only a few cars. Although the environmental light frequently changes (an example is the lilac section on the left on the Original Timeline), the count of false detections is low. Fused cars in the middle are split appropriately. However the stopping pedestrian (top right corner) is ignored, al well as the other two pedestrians who are falsely split based on their shapes.}
\end{sidewaysfigure}

\centering
\begin{sidewaysfigure}[h!]
	\centering
	\captionsetup{width=0.9\textwidth,aboveskip=10pt}
	\includegraphics[width=\textwidth]{ErC1.png}
	\caption[Some output Timelines of the ErC1 test video]{The figures (from the top to the bottom) are as follows: Original Timeline Image, Result Timeline Image and Evaluation Result Timeline Image respectively. The video captures a stopping tram, a handful of cars and pedestrians. The tram's slow movement is still a challenge to the fusing algorithm, thus causes some noise.}%{ErC kódú videó futtatási eredménye. Fentről lefelé a képek rendre a következők. entről lefelé a képek rendre a következők. Eredeti videóból képzett TK. MOG2 TK. A kiértékelés eredménye. \\
		%A videón egy megálló villamos, illetve a mellette haladó gyalogosok és autók láthatók. A rendszer jól összeilleszti a megálló villamos blobját, de helytelenül a jobb alsó sarokban található blobokat (hosszú téglalap), amely csak a a háttér folytonossága miatt tűnik a rendszernek megálló autónak.}
\end{sidewaysfigure}

\centering
\begin{sidewaysfigure}[h!]
	\centering
	\captionsetup{width=0.9\textwidth,aboveskip=10pt}
	\includegraphics[width=\textwidth]{SasA1.png}
	\caption[Some output Timelines of the SasA1 test video]{The figures (from the top to the bottom) are as follows: Original Timeline Image, MOG2 Timeline Image, Speed Timeline Image and Evaluation Result Timeline Image respectively. \\
	The software performs well on this video. Most fused vehicles are correctly split, and stopping vehicles' blobs are reattached as well. Due to shaky camera, that causes the background to change, some false alarms are present (on the top, middle part). The Speed Timeline is calculated only in the centre of road lanes displaying fewer values than the Parameter Timelines in figures \ref{fig:ErE5}., \ref{fig:GoldA1}., \ref{fig:GoldB1}.}
\end{sidewaysfigure}




\label{page:last}
\end{document}
